\section{辛流形的自由环路空间}
\label{sec:free-loop-space-symplectic-manifold}

现在把上一节的讨论专门化到辛流形 $(M,\omega)$ 的情形。

\subsection{\texorpdfstring{$S^1$-作用及其矩映射}{S¹-作用及其矩映射}}
\label{subsec:circle-action-moment-map}

此时 $\mathcal L(M)$ 带有如下定义的典范（弱）辛形式 $\Omega$：
\begin{equation}
 \Omega(\xi_1,\xi_2):=\int_0^1\omega(\xi_1(t),\xi_2(t))\,dt.
 \label{eq:18-loop-space-symplectic-form}
\end{equation}
我们概述这个 2-形式的一些基本众所周知性质，而不试图给出完全严格的证明。关于一般
环路空间上的微分形式，可参见 (Chen73)、(GJP91) 的严谨论述。

\begin{proposition}
\label{prop:18-loop-space-symplectic-properties}
\begin{enumerate}[label=(\arabic*)]
\item 形式 $\Omega$ 是弱非退化的。换言之，对任意 $\gamma\in\mathcal L(M)$，
若对所有 $\eta\in T_\gamma\mathcal L(M)$ 都有 $\Omega(\xi,\eta)=0$，则且仅则
$\xi=0$。
\item $\Omega$ 是闭的。
\item 式~\eqref{eq:18-loop-rotation} 的 $S^1$-作用是辛的，即保持
$\Omega$；也就是说 $\mathcal L_{\mathbb X}\Omega=0$。
\footnote{译者注：原书此项引用式 (18.1.5)，但该式是覆盖空间上二元组
$(\gamma,w)$ 的提升作用；本命题讨论 $\mathcal L(M)$ 上的 $\Omega$，相应的定义域
旋转是式 (18.1.3)，译文据此改正引用。}
\end{enumerate}
\end{proposition}

$\Omega$ 的（弱）非退化性来自 $\omega$ 的非退化性；$\Omega$ 的闭性来自
$\omega$ 的闭性以及 \emph{$S^1$ 没有边界这一事实}。有关计算留给读者，也可参见 (Wn78)
对这些陈述的验证。

于是，通过投影 $\widetilde{\mathcal L}_0(M)\to\mathcal L_0(M)$ 的拉回，形式
$\Omega$ 在覆盖空间 $\widetilde{\mathcal L}_0(M)$ 上诱导一个辛形式，记为
$\widetilde\Omega$。

\stfhIndex{I050}{Dianfanbozhumian}{典范薄柱面（canonical thin cylinder）}
为说明上面陈述 (3) 的含义，先回忆两个邻近环路之间的\emph{典范薄柱面}这一有用
概念。记 $\exp$ 为度量
\[
 g:=\omega(\,\cdot,J_{\mathrm{ref}}\,\cdot)
\]
的指数映射，其中 $J_{\mathrm{ref}}$ 是任意给定的固定相容殆复结构。记
$\iota(g)$ 为 $g$ 的\termInjectivityRadius{}。只要给定的两点满足
$d(x,y)<\iota(g)$，就可以唯一地写成
\[
 y=\exp_x(E(x,y)),\qquad E(x,y):=(\exp_x)^{-1}(y).
\]
因此，如果环路 $\gamma,\gamma':S^1\to M$ 之间的 $C^0$ 距离
$d_{C^0}(\gamma,\gamma')$ 小于 $\iota(g)$，就可以定义典范映射
\[
 u^{\mathrm{can}}_{\gamma\gamma'}:[0,1]\times S^1\longrightarrow M
\]
为
\begin{equation}
 u^{\mathrm{can}}_{\gamma\gamma'}(s,t)
 =\exp_{\gamma(t)}\bigl(sE(\gamma,\gamma')(t)\bigr),
 \label{eq:18-canonical-thin-cylinder}
\end{equation}
其中 $E(\gamma,\gamma')(t):=(\exp_{\gamma(t)})^{-1}(\gamma'(t))$。

\begin{lemma}
\label{lem:18-action-one-form-locally-exact}
设 $\mathbb X$ 是式~\eqref{eq:18-circle-generator} 中的向量场。形式
$\mathbb X\mathbin{\lrcorner}\Omega$ 是 $\mathcal L(M)$ 上的闭 1-形式；具体说，
在任意给定环路 $\gamma_0$ 的一个 $C^1$ 邻域上，都存在函数
$\mathcal A=\mathcal A_0$，满足
\begin{equation}
 d\mathcal A=\mathbb X\mathbin{\lrcorner}\Omega.
 \label{eq:18-local-action-primitive}
\end{equation}
\end{lemma}

\begin{proof}
对任意与给定 $\gamma_0$ 在 $C^\infty$ 意义下充分接近的环路 $\gamma$，考虑
\[
 u^{\mathrm{can}}_{\gamma_0\gamma}:s\in[0,1]
 \longmapsto\exp_{\gamma_0}\bigl(sE(\gamma_0,\gamma)\bigr),
\]
它定义了端点固定的道路的一个特定同伦类
$[u_{\gamma_0\gamma}]$，且 $u(0)=\gamma_0$、$u(1)=\gamma$。相伴的参数化柱面
\[
 [0,1]\times S^1\longrightarrow M,\qquad
 (s,t)\longmapsto\exp_{\gamma_0(t)}\bigl(sE(\gamma_0(t),\gamma(t))\bigr)
\]
也记作 $u^{\mathrm{can}}_{\gamma_0\gamma}$。从它的显式表达式可知，在
$\gamma_0\in\mathcal L(M)$ 的一个 $C^1$-小邻域中，可以用公式
\begin{equation}
 \mathcal A(\gamma;\gamma_0)=0-\int
   (u^{\mathrm{can}}_{\gamma_0\gamma})^*\omega
 \label{eq:18-local-action-function}
\end{equation}
定义函数 $\mathcal A$。这里的“$0$”应理解为可任意选取的值
$\mathcal A(\gamma_0;\gamma_0)$。

下面验证~\eqref{eq:18-local-action-primitive}。设 $\gamma$ 位于给定的
$C^1$-小邻域中，$\xi\in T_\gamma\mathcal L(M)$。用道路芽
$\{\gamma_s\}_{-\epsilon<s<\epsilon}$ 表示 $\xi$，并对
$-\epsilon<s<\epsilon$ 考虑二元组 $(\gamma_s,u_s)$，其中
\[
 u_s:=u^{\mathrm{can}}_{\gamma_0\gamma}\#\{\gamma_s\}_{[0,s]}.
\]
计算得
\begin{align}
 d\mathcal A(\gamma)(\xi)
 &=\left.\frac d{ds}\right|_{s=0}\left(-\int u_s^*\omega\right)\notag\\
 &=\left.\frac d{ds}\right|_{s=0}\left(
   -\int (u^{\mathrm{can}}_{\gamma_0\gamma})^*\omega
   -\int_0^s\int_0^1
   \omega\left(\frac{\partial\gamma_u(t)}{\partial u},
                \frac{\partial\gamma_u(t)}{\partial t}\right)dt\,du
   \right)\notag\\
 &=\int_0^1-\omega(\xi(t),\dot\gamma(t))\,dt
  =(\mathbb X(\gamma)\mathbin{\lrcorner}\Omega)(\xi).
 \label{eq:18-action-one-form-computation}
\end{align}
这与~\eqref{eq:18-action-one-form-computation} 一起证明了
\eqref{eq:18-local-action-primitive}，从而证明引理。
\end{proof}

\begin{remark}
\label{rem:18-loop-space-de-rham-class}
从环路空间德拉姆理论 (Chen73)、(GJP91) 的观点看，$M$ 上的辛形式 $\omega$
通过闭 1-形式 $\mathbb X\mathbin{\lrcorner}\Omega$ 诱导一个典范的 1 次上同调类。
\stfhIndex{I115}{Diedaijifen}{迭代积分（iterated integral）}
用 (GJP91) 第 344 页的记号，该类由\emph{迭代积分} $\widetilde P_1(\omega)$ 得到。这个
\stfhIndex{I196}{Ruozhenghe}{弱正合（weakly exact）}
1-形式一般\emph{不是正合}的；它正合，恰好就是辛形式 $\omega$ 所谓的\emph{弱正合性}。
\end{remark}

若把这个闭 1-形式限制到 $\mathcal L_0(M)$，再考虑它在万有覆盖空间
$\widetilde{\mathcal L}_0(M)$ 上的提升，则公式~\eqref{eq:18-local-action-function}
可由二元组 $(\gamma,w)$ 的函数推广为\emph{全局}提升，仍记作
$\mathcal A=\mathcal A_0$：
\[
 \mathcal A_0(\gamma,w)=-\int w^*\omega.
\]
这里把 $w$ 看作从常道路 $w(0)$ 到 $\gamma=\partial w$ 的道路。积分前取负号，是为了
与作者关于谱不变量的一系列论文（例如 (Oh06a)）所用约定一致。称
$\mathcal A_0$ 为\emph{未扰动作用量泛函}。它在
$\widetilde{\mathcal L}_0(M)$ 上满足
\begin{equation}
 d\mathcal A_0=\mathbb X\mathbin{\lrcorner}\Omega.
 \label{eq:18-global-action-primitive}
\end{equation}
换言之，$\widetilde{\mathcal L}_0(M)$ 上的 $S^1$-作用是哈密顿的，与之相伴的矩映射
正是函数 $\mathcal A_0:\widetilde{\mathcal L}_0(M)\to\mathbb R$。更详细的讨论见
(Wn78)。

\subsection{\personNovikov{}覆盖}
\label{subsec:novikov-covering}

万有覆盖空间虽提供了使作用 1-形式的提升成为正合形式的模型，却太大而无法在其上
开展有意义的临界点理论：相伴函数的对称群
$\pi_1(\mathcal L_0(M))\cong\pi_2(M)\times\pi_1(M)$（参见
式~\eqref{eq:18-free-loop-fundamental-group} 后的译者注）可能高度\emph{非交换}，除非 $M$
单连通。因此要寻找一个更小的覆盖空间，使其覆盖变换群较小且为交换群。\personNovikov{}
(No81) 最先通过考虑一个使闭 1-形式的拉回成为正合形式的\emph{循环}覆盖，发展了紧流形上
闭 1-形式的\personMorse{}理论，并用 $M$ 的拓扑数据估计闭 1-形式的零点数。

\stfhIndex{I138}{Nuoweikefufugai}{诺维科夫覆盖（Novikov covering）}
依照 (Fl89b)、(HS95)，现在引入 $\mathcal L_0(M)$ 的\emph{\termNovikovCovering{}}概念。
它是 $\mathcal L_0(M)$ 的一个交换覆盖空间，对应于\personNovikov{} \personMorse{}理论中的循环覆盖。
沿用\personSeidel{} (Se97) 的术语，作如下定义。

\begin{definition}
\label{def:18-gamma-equivalence}
设 $(\gamma,w)$ 是由 $\gamma\in\mathcal L_0(M)$ 和以 $\gamma$ 为边界的圆盘 $w$
\stfhIndex{I007}{00-gamma-equivalence}{$\Gamma$-等价关系（$\Gamma$-equivalence）}
组成的二元组。称 $(\gamma,w)$ 与 $(\gamma,w')$ \emph{$\Gamma$-等价}，当且仅当
\[
 \omega([w'\#\overline w])=0,\qquad
 c_1([w'\#\overline w])=0,
\]
其中 $\overline w$ 表示在定义域上取相反定向的映射，$w'\#\overline w$ 是显然的
粘合球面。这里 $\Gamma$ 表示群
\[
 \Gamma=\frac{\pi_2(M)}
 {\ker(\omega|_{\pi_2(M)})\cap\ker(c_1|_{\pi_2(M)})}.
\]
记
\[
 \Gamma_\omega:=\omega(\Gamma)=\omega(\pi_2(M))\subset\mathbb R,
\]
\stfhIndex{I151}{Zhouqiqun}{周期群（period group）}
并称它为 $(M,\omega)$ 的（球面）\emph{\termPeriodGroup{}}。
\end{definition}

\begin{definition}
\label{def:18-rational-symplectic-manifold}
\stfhSubIndex{I186}{Xinliuxing}{辛流形（symplectic manifold）}{Youli}{有理（rational）}
若 $\Gamma_\omega\subset\mathbb R$ 是离散子群，则称 $(M,\omega)$ 为\emph{有理的}；
\stfhSubIndex{I185}{Xinliuxing}{辛流形（symplectic manifold）}{Wuli}{无理（irrational）}
否则称其为\emph{无理的}。
\end{definition}

\begin{example}
\label{ex:18-product-spheres-rationality}
带有乘积辛形式 $\omega_1\oplus\omega_2$ 的乘积
$S^2(r_1)\times S^2(r_2)$ 是有理的，当且仅当比值 $r_2^2/r_1^2$ 是有理数。
\end{example}

\begin{remark}
\label{rem:18-irrational-period-group}
注意，对无理的 $(M,\omega)$，周期群是 $\mathbb R$ 的可数稠密子集。一般认为，无理
辛流形上的哈密顿流，其动力学行为会比有理辛流形上的情形复杂得多。周期群
$\Gamma_\omega$ 是两者这种不同动力学行为的最简单指标。
\end{remark}

从现在起，专用 $[\gamma,w]$ 表示 $(\gamma,w)$ 的 $\Gamma$-等价类，并以
$\widetilde{\mathcal L}_0(M)$ 表示这些等价类的空间。记典范投影为
$\pi:\widetilde{\mathcal L}_0(M)\to\mathcal L_0(M)$，称
$\widetilde{\mathcal L}_0(M)$ 为 $\mathcal L_0(M)$ 的\termGammaCoveringSpace{}。
对 $A\in\pi_2(M)$，用 $A$ 或 $q^A$ 表示它在投影
$\pi_2(M)\to\Gamma$ 下的像。与 $A$ 相伴有两个自然不变量：\emph{赋值} $\nu(A)$，
\begin{equation}
 \nu:\Gamma\longrightarrow\mathbb R,\qquad \nu(A)=\omega(A),
 \label{eq:18-gamma-valuation}
\end{equation}
以及\emph{次数} $d(A)$，
\begin{equation}
 d:\Gamma\longrightarrow\mathbb Z,\qquad d(A)=2c_1(A).
 \label{eq:18-gamma-degree}
\end{equation}
一般而言，这两个不变量相互独立（见 (Gom98)），所以 $q^A$ 是依赖于两个变量的
形式参数。从这个意义说，也可以写成
\[
 q^A=T^{\omega(A)}e^{c_1(A)},
\]
其中 $T$ 与 $e$ 是两个不同的形式参数。规定 $T$ 的次数为 $0$，$e$ 的次数为 $2$。
\footnote{译者注：原书把 $e$ 的指数印作 $c_1(A)/2$，同时规定
$d(A)=2c_1(A)$ 且 $e$ 的次数为 $2$，两者给出的 $q^A$ 次数不一致；下一章又明确使用
$q^A=T^{\omega(A)}e^{c_1(A)}$。译文按次数定义及下一章记号补正。}

上面定义的（未扰动）作用量泛函 $\mathcal A_0$ 显然可下降到 $\Gamma$-覆盖空间，
公式仍为
\[
 \mathcal A_0([\gamma,w])=-\int w^*\omega.
\]
这个泛函在 $\widetilde{\mathcal L}_0(M)$ 上给出一个自然递增滤过：对每个
$\lambda\in\mathbb R$，定义
\[
 \widetilde{\mathcal L}_0^\lambda(M):=
 \{[z,w]\in\widetilde{\mathcal L}_0(M)\mid
   \mathcal A_0([z,w])\leq\lambda\}.
\]
注意，当 $\lambda\leq\lambda'$ 时，
\[
 \widetilde{\mathcal L}_0^\lambda(M)
 \subset\widetilde{\mathcal L}_0^{\lambda'}(M).
\]
由~\eqref{eq:18-global-action-primitive}，$\mathcal A_0$ 的临界集
$\operatorname{Crit}\mathcal A_0$ 是若干份 $M$ 的不交并：
\[
 \operatorname{Crit}\mathcal A_0=\bigsqcup_{g\in\Gamma}g\cdot M,
\]
其中 $M\hookrightarrow\mathcal L_0(M)$，$x\mapsto[x,\widetilde x]$ 是典范嵌入，
$[x,\widetilde x]$ 是常环路 $x$ 与常圆盘 $\widetilde x$ 组成的二元组。我们有自然交换图
\[
 \begin{array}{ccc}
  \operatorname{Crit}\mathcal A_0&\longrightarrow&
    \widetilde{\mathcal L}_0(M)\\
  \big\downarrow\scriptstyle\pi&&\big\downarrow\scriptstyle\pi\\
  M&\longrightarrow&\mathcal L_0(M),
 \end{array}
\]
其公共纤维同构于 $\Gamma$。

\subsection{第二变分}
\label{subsec:unperturbed-action-second-variation}

现在计算 $\mathcal A_0$ 在每个临界点 $[x,\widetilde x\#A]$ 处的第二变分
$d^2\mathcal A_0$。利用覆盖投影的切映射所诱导的典范认同
\begin{equation}
 T_{[x,\widetilde x]}\widetilde{\mathcal L}(M)
 \cong T_x\mathcal L(M),
 \label{eq:18-cover-tangent-identification}
\end{equation}
并把 $d^2\mathcal A_0([x,\widetilde x\#A])$ 看作定义在
$T_x\mathcal L(M)$ 上的二次型。

\begin{proposition}
\label{prop:18-unperturbed-action-bott-morse}
在每个 $[x,\widetilde x\#A]\in\operatorname{Crit}\mathcal A_0$ 处，Hessian
$d^2\mathcal A_0$ 在
\[
 T_{[x,\widetilde x\#A]}\widetilde{\mathcal L}_0(M)
 \cong T_x\mathcal L_0(M)
\]
上定义一个双线性型；它在 $\mathcal L_0(M)$ 中
$\operatorname{Crit}\mathcal A_0$ 的法方向上（弱）非退化。
\end{proposition}

\begin{proof}
为简化记号，记
\[
 [x,\widetilde x]=\widetilde x,\qquad
 [x,\widetilde x\#A]=\widetilde x\otimes q^A.
\]
由于作用 $\#A$ 是覆盖空间
$\widetilde{\mathcal L}_0(M)\to\mathcal L_0(M)$ 的覆盖变换，只需考虑点
$[x,\widetilde x]$。设
$V:(-\epsilon,\epsilon)\times(-\epsilon,\epsilon)\to\mathcal L(M)$ 是二参数环路族，
满足
\[
 \begin{cases}
 V(0,0,t)\equiv x,\\
 \left.\dfrac{\partial V}{\partial s}\right|_{s=u=0}(t)=\xi_1(t),\quad
 \left.\dfrac{\partial V}{\partial u}\right|_{s=u=0}(t)=\xi_2(t).
 \end{cases}
\]
以 $\gamma_{u,s}:S^1\to M$ 表示环路
$\gamma_{u,s}(t):=V(u,s,t)$，并令 $w_{u,s}:D^2\to M$ 为满足
$w_{u,s}|_{\partial D^2}=\gamma_{u,s}$ 的“小”圆盘。这样的圆盘在同伦意义下唯一。

由第二变分的定义，
\[
 d^2\mathcal A_0([x,\widetilde x])(\xi_1,\xi_2)
 =\left.\frac{\partial^2}{\partial u\,\partial s}\right|_{s=u=0}
  \mathcal A_0([\gamma_{u,s},w_{u,s}]).
\]
这里通过认同~\eqref{eq:18-cover-tangent-identification}，把 $\xi_i$ 看成
$T_x\mathcal L(M)$ 中的切向量。由~\eqref{eq:18-global-action-primitive}，
\[
 \frac\partial{\partial s}\mathcal A_0([\gamma_{u,s},w_{u,s}])
 =-\int\omega\left(\frac{\partial V}{\partial t},
                     \frac{\partial V}{\partial s}\right)dt.
\]
因此，使用任意无挠联络 $\nabla$，得到
\begin{align*}
 &\left.\frac{\partial^2}{\partial u\,\partial s}\right|_{s=u=0}
 \mathcal A_0([\gamma_{u,s},w_{u,s}])\\
 &=-\int_0^1(\nabla_u\omega)(x)
   \left(\left.\frac{\partial V}{\partial t}\right|_{u=s=0},
         \left.\frac{\partial V}{\partial s}\right|_{u=s=0}\right)dt\\
 &\quad-\int_0^1\omega\left(
   \left.\nabla_u\frac{\partial V}{\partial t}\right|_{u=s=0},
   \left.\frac{\partial V}{\partial s}\right|_{u=s=0}\right)dt
 -\int_0^1\omega\left(
   \left.\frac{\partial V}{\partial t}\right|_{u=s=0},
   \left.\nabla_u\frac{\partial V}{\partial s}\right|_{u=s=0}\right)dt.
\end{align*}
因为 $V(0,0,t)\equiv x$，故
$\left.\partial V/\partial t\right|_{u=s=0}\equiv0$，第一项和第三项均消失。对第二项，
由 $\nabla$ 无挠可得
\[
 \nabla_u\frac{\partial V}{\partial t}
 =\nabla_t\frac{\partial V}{\partial u}.
\]
从而
\[
 \left.\nabla_u\frac{\partial V}{\partial t}\right|_{u=s=0}
 =\left.\nabla_t\frac{\partial V}{\partial u}\right|_{u=s=0}
 =\left.\nabla_t\xi_2\right|_{u=s=0}
 =\left.\frac{d\xi_2}{dt}(t)\right|_{u=s=0},
\]
最后一个等式是因为 $\xi_2$ 是沿常环路 $x$ 的变分。综合可得
\[
 d^2\mathcal A_0([x,\widetilde x])(\xi_1,\xi_2)
 =-\int_0^1\omega\left(\frac{d\xi_2}{dt}(t),\xi_1(t)\right)dt.
\]
特别地，$d^2\mathcal A_0([x,\widetilde x])$ 的核由沿常道路 $\widetilde x$ 的向量场
\[
 \left\{\xi_2\in C^\infty(S^1,T_xM)\ \middle|\
 \frac{d\xi_2}{dt}=0\right\}
\]
组成。通过嵌入 $M\hookrightarrow C^\infty(S^1,M)$，可以把 $T_xM$ 与
$C^\infty(S^1,T_xM)$ 中的像 $T_xM$ 认同。因此
\[
 \ker d^2\mathcal A_0([x,\widetilde x])\cong T_xM,
\]
也就是说，该核恰好是 $\mathcal A_0$ 的临界流形
$\operatorname{Crit}\mathcal A_0$ 在任意点
$[x,\widetilde x]\in\operatorname{Crit}\mathcal A_0$ 处的切空间。对任意
$A\in\Gamma$，其他临界点 $[x,\widetilde x\#A]$ 也同样成立。
\end{proof}

按此命题，$\mathcal A_0$ 可视为\personBott{}--\personMorse{}函数 (Bo54)。我们将看到，加入通有的
含时哈密顿函数 $H:[0,1]\times M\to\mathbb R$ 后，可以消去这些临界点中的绝大
多数，只留下 $H$ 的有限个\emph{非退化} 1-周期哈密顿轨道。\personArnold{}的著名猜想 (Ar65)
断言，相伴哈密顿流至少有 $SB(M)$（即 $H^*(M)$ 的秩）个（可缩）1-周期轨道。
